\section{Results}
\label{sec:results}

\subsection{Measured telemetry}

Before any model was loaded, the workers reported the figures in
Table~\ref{tab:telemetry}. These are the values the scheduler consumes.

\begin{table}[htbp]
\centering\small
\caption{Worker telemetry as received by the controller. Dashes mark fields the
platform does not expose; these are transmitted as \code{null} and shown in the
dashboard as \textit{Unavailable}.}
\label{tab:telemetry}
\begin{tabular}{@{}lrr@{}}
\toprule
 & \textbf{TUF (Windows)} & \textbf{M4 Mac} \\
\midrule
GPU name           & RTX 4050 Laptop   & Apple M4 \\
GPU memory, free / total & \SI{5843}{} / \SI{6141}{MiB} & --- \\
GPU utilisation    & \SI{0}{\percent}  & --- \\
GPU temperature    & \SI{50}{\celsius} & --- \\
GPU power draw     & \SI{4.09}{W}      & --- \\
GPU power limit    & ---               & --- \\
System RAM, available / total & \SI{4845}{} / \SI{16003}{MiB} & \SI{5938}{} / \SI{16384}{MiB} \\
Battery, AC        & \SI{84}{\percent}, online & \SI{16}{\percent}, online \\
\bottomrule
\end{tabular}
\end{table}

The Windows GPU power limit is blank because \code{nvidia-smi} returned
\code{[N/A]} on that machine. The Mac's GPU rows are blank because macOS
publishes no per-GPU memory, utilisation, temperature or power figure for Apple
Silicon through the interfaces we use. Both cases are reported as unavailable
rather than filled with a substitute.

\subsection{Scheduling and actual placement}

With both NVIDIA devices reporting free memory, the scheduler produced weights
directly from measurement:

\[
\underbrace{5843 - 512}_{\text{TUF RPC}} = 5331,
\qquad
\underbrace{3716 - 512}_{\text{ROG local}} = 3204,
\qquad\Rightarrow\quad
\texttt{--tensor-split 5331,3204}
\]

Figure~\ref{fig:placement} shows the placement \texttt{llama.cpp} subsequently
reported for four configurations. These are read back from its load output, not
inferred from the request.

\begin{figure}[htbp]
\centering
\begin{tikzpicture}[font=\scriptsize]
\begin{axis}[hpaxis,
  width=11.6cm,height=5.0cm,
  ybar stacked, bar width=18pt,
  ymin=0,ymax=5200,
  ylabel={model weights (MiB)},
  symbolic x coords={3B three machines,3B after Mac loss,3B after TUF loss,8B two machines},
  xtick=data,
  x tick label style={font=\tiny,align=center},
  legend style={at={(0.5,1.14)},anchor=north,legend columns=4,draw=gray!45},
  nodes near coords={},
  grid=major,
  ]
% CPU-mapped
\addplot+[fill=serC!35,draw=serC!70,bar shift=0pt]
  coordinates {(3B three machines,308) (3B after Mac loss,308) (3B after TUF loss,308) (8B two machines,282)};
% ROG local CUDA
\addplot+[fill=serA!55,draw=serA,bar shift=0pt]
  coordinates {(3B three machines,552) (3B after Mac loss,896) (3B after TUF loss,1918) (8B two machines,1073)};
% TUF RPC
\addplot+[fill=green!45!black!55,draw=green!35!black,bar shift=0pt]
  coordinates {(3B three machines,398) (3B after Mac loss,1022) (3B after TUF loss,0) (8B two machines,3330)};
% Mac RPC
\addplot+[fill=orange!65,draw=orange!80!black,bar shift=0pt]
  coordinates {(3B three machines,968) (3B after Mac loss,0) (3B after TUF loss,0) (8B two machines,0)};
\legend{CPU-mapped (ROG RAM), ROG CUDA0, TUF RPC0, Mac RPC0}
\end{axis}
\end{tikzpicture}
\caption{Per-device model weights reported by \texttt{llama.cpp} after loading,
for four measured configurations. The first two bars are the same request
before and after the Mac accelerator was stopped: its \SI{968}{MiB} is
redistributed onto the two surviving machines. The fourth bar is the 8B
capacity run, where the remote Windows GPU holds \SI{3330}{MiB} of the
\SI{4.9}{GB} model.}
\label{fig:placement}
\end{figure}

Two observations follow from the figure. In the three-machine configuration the
Mac received the largest remote share, \SI{968}{MiB}, even though HyperPod
supplied no weights for it; that allocation was \texttt{llama.cpp}'s own
decision under the documented fallback. And in the 8B run the remote Windows GPU
holds \SI{3330}{MiB} against \SI{1073}{MiB} on the controller. That is
\SI{71}{\percent} of all reported model weights, or \SI{76}{\percent} of the
weights resident on a GPU, and it is what makes that model runnable at all on
this hardware.

\subsection{Latency and throughput}
\label{sec:perf}

Figure~\ref{fig:perf} gives all five runs for each configuration. We plot every
run rather than only the medians, because the run-to-run spread is the most
important feature of the distributed measurements.

\begin{figure}[htbp]
\centering
\begin{tikzpicture}[font=\scriptsize]
\begin{groupplot}[group style={group size=2 by 1,horizontal sep=1.45cm},
  hpaxis,width=6.6cm,height=5.0cm,xmin=0.6,xmax=5.4,xtick={1,2,3,4,5},
  xlabel={run},]

\nextgroupplot[title={End-to-end latency},ylabel={seconds},ymin=0,ymax=17.5,
  legend style={at={(1.28,1.40)},anchor=north,legend columns=3,draw=gray!45}]
\addplot+[dottedC] coordinates {(1,1.566)(2,1.527)(3,1.623)(4,1.543)(5,1.505)};
\addlegendentry{local baseline}
\addplot+[dashedB] coordinates {(1,7.122)(2,15.025)(3,8.157)(4,13.807)(5,11.947)};
\addlegendentry{native split}
\addplot+[solidA] coordinates {(1,10.739)(2,10.920)(3,15.930)(4,11.804)(5,14.324)};
\addlegendentry{HyperPod split}

\nextgroupplot[title={Decode throughput},ylabel={tokens / second},ymin=0,ymax=85]
\addplot+[dottedC] coordinates {(1,75.64)(2,75.68)(3,75.51)(4,75.54)(5,75.45)};
\addplot+[dashedB] coordinates {(1,18.13)(2,9.17)(3,20.91)(4,9.94)(5,13.19)};
\addplot+[solidA] coordinates {(1,15.30)(2,19.04)(3,9.12)(4,11.79)(5,9.94)};
\end{groupplot}
\end{tikzpicture}
\caption{Five sequential requests per configuration against a warm model.
The local baseline is tightly clustered, varying by \SI{0.12}{\second} and
\SI{0.23}{tokens/second} across all five runs. Both distributed configurations
vary by several seconds and by more than a factor of two in throughput, which
we attribute to the Wi-Fi link rather than to the control plane.}
\label{fig:perf}
\end{figure}

\begin{table}[htbp]
\centering\small
\caption{Medians and ranges over the five runs of Figure~\ref{fig:perf}.}
\label{tab:perf}
\begin{tabular}{@{}lrrrr@{}}
\toprule
\textbf{Configuration} & \textbf{Median latency} & \textbf{Range} & \textbf{Median decode} & \textbf{Range} \\
\midrule
Local baseline  & \SI{1.543}{\second}  & 1.505--1.623 & \SI{75.54}{} & 75.45--75.68 \\
Native split    & \SI{11.947}{\second} & 7.122--15.025 & \SI{13.19}{} & 9.17--20.91 \\
HyperPod split  & \SI{11.804}{\second} & 10.739--15.930 & \SI{11.79}{} & 9.12--19.04 \\
\bottomrule
\end{tabular}
\end{table}

The local baseline decoded at a median \SI{75.54}{tokens/second} against
\SI{11.79}{} for the HyperPod split: distributing this model over this Wi-Fi
network cost roughly a factor of six in throughput. That is the central cost of
the approach and we do not present it as anything else.

Between the two distributed configurations the median latencies differ by
\SI{0.143}{\second}, which is far smaller than the spread within either
configuration: \SI{5.2}{\second} for the HyperPod split and \SI{7.9}{\second}
for the native split. With five runs we can state that the control plane
introduced no latency penalty large enough to separate from link variance; we
cannot make a stronger claim than that, and a larger sample would be needed to
try.

\subsection{Failure and recovery}
\label{sec:recovery}

Three failure experiments were run, all by stopping a real
\code{ggml-rpc-server} process by hand during a live request.
Figure~\ref{fig:timeline} shows them on a common time axis.

\begin{figure}[htbp]
\centering
\begin{tikzpicture}[font=\scriptsize,
  lbl/.style={anchor=east,font=\tiny},
  bar/.style={line width=8pt},
  tick/.style={gray!60,line width=0.4pt},
  outlbl/.style={font=\tiny,anchor=west}]

\def\sx{0.135}   % seconds -> cm

% legend
\begin{scope}[shift={(0,2.95)}]
  \foreach \i/\c/\t in {0/serA!45/generating, 1/serB!55/abort interval,
                        2/orange!50/model reload, 3/green!45!black!55/retried request}{
    \draw[bar,\c] (\i*2.55,0) -- (\i*2.55+0.38,0);
    \node[font=\tiny,anchor=west] at (\i*2.55+0.46,0) {\t};
  }
\end{scope}

% axis
\draw[tick] (0,-0.5) -- (80*\sx,-0.5);
\foreach \t in {0,10,20,30,40,50,60,70,80}{
  \draw[tick] (\t*\sx,-0.5) -- (\t*\sx,-0.63) node[below,font=\tiny]{\t};
}
\node[font=\tiny] at (40*\sx,-1.1) {seconds after the request was submitted};

% --- native llama.cpp ---
\node[lbl] at (-0.12,2.1) {native \texttt{llama.cpp}};
\draw[bar,serA!45] (0,2.1) -- (15*\sx,2.1);
\draw[bar,serB!55] (15*\sx,2.1) -- (34.06*\sx,2.1);
\draw[serB,line width=0.6pt] (34.06*\sx,1.92) -- (34.06*\sx,2.28);
\node[outlbl,text=serB] at (34.9*\sx,2.1) {request failed, HTTP 000};

% --- HyperPod, Mac loss ---
\node[lbl] at (-0.12,1.3) {HyperPod, Mac stopped};
\draw[bar,serA!45] (0,1.3) -- (21*\sx,1.3);
\draw[bar,serB!55] (21*\sx,1.3) -- (22*\sx,1.3);
\draw[bar,orange!50] (22*\sx,1.3) -- (54*\sx,1.3);
\draw[bar,green!45!black!55] (54*\sx,1.3) -- (75.68*\sx,1.3);
\draw[green!35!black,line width=0.6pt] (75.68*\sx,1.12) -- (75.68*\sx,1.48);
\node[outlbl,text=green!35!black] at (76.5*\sx,1.3) {returned};

% --- HyperPod, TUF loss ---
\node[lbl] at (-0.12,0.5) {HyperPod, TUF stopped};
\draw[bar,serA!45] (0,0.5) -- (4*\sx,0.5);
\draw[bar,serB!55] (4*\sx,0.5) -- (6*\sx,0.5);
\draw[bar,orange!50] (6*\sx,0.5) -- (7*\sx,0.5);
\draw[bar,green!45!black!55] (7*\sx,0.5) -- (11.31*\sx,0.5);
\draw[green!35!black,line width=0.6pt] (11.31*\sx,0.32) -- (11.31*\sx,0.68);
\node[outlbl,text=green!35!black] at (12.1*\sx,0.5) {returned after \SI{11.31}{\second}};
\end{tikzpicture}
\caption{The three physical failure experiments. Blue is generation, red is the
interval containing the abort, orange is model reload, green is the retried
request. In the top run, unmodified \texttt{llama.cpp} never produced a
response. In both HyperPod runs the original HTTP request completed without
resubmission. Prompt lengths differed between runs, so the total durations are
not comparable to each other; the comparison being made is outcome, not speed.
The moment the operator stopped the native run's accelerator is known only to
the nearest minute, so its red interval begins at an approximate offset; all
other boundaries come from logged timestamps.}
\label{fig:timeline}
\end{figure}

\paragraph{Native baseline.}
A request was submitted against the two-machine split and the Windows
accelerator was stopped roughly fifteen seconds later. The client received no
response: \code{curl} exited 52 with HTTP~000 after \SI{34.06}{\second}. The
model process had aborted and did not return.

\paragraph{Three machines to two.}
The 3B model was loaded across all three machines and a request capped at 420
output tokens was submitted. The Mac's \code{ggml-rpc-server} was then stopped,
while its Rust agent was deliberately left running. This separates two states
that are easy to conflate: the agent continued to heartbeat normally, so agent
health alone would have indicated a working node.

The controller log records the sequence:

\begin{quote}\small\ttfamily
llama-server gen 1 exited: signal: aborted\\
infer attempt 1 recognized node death and re-ran the same request\\
pool skipping 192.168.1.13:50052 (worker heartbeat ok, but its rpc-server is unreachable)\\
launched gen 2 over [192.168.1.5:50052]; scheduler=measured-free-vram weights=4763,3204\\
pipeline ready (gen 2) after 33s
\end{quote}

The rebuilt pipeline was still distributed, now holding \SI{896}{MiB} on the
controller GPU and \SI{1022}{MiB} on the Windows worker. The original request
returned a 1{,}937-character response \SI{75.68}{\second} after submission.

\paragraph{Two machines to one.}
Repeating the experiment against the Windows worker produced the same outcome by
a shorter path. Generation~1 aborted four seconds after submission, the
unreachable endpoint was excluded, and because the 3B model fits in the
controller's GPU alone the rebuild completed in about one second. The same
request returned a 1{,}828-character response after \SI{11.31}{\second} in
total.

\subsection{Load times}
\label{sec:load}

Loading dominates the time cost of this design, and it varies with cache state
and link quality rather than with anything HyperPod controls.
Figure~\ref{fig:load} collects the measured loads.

\begin{figure}[htbp]
\centering
\begin{tikzpicture}[font=\scriptsize]
\begin{axis}[hpaxis,
  width=12.2cm,height=4.4cm,
  xbar, bar width=11pt,
  xmin=0,xmax=1080,
  xlabel={seconds to a ready model},
  symbolic y coords={{3B rebuild, local only},{3B rebuild, 2 machines},{3B cached, 3 machines},{3B cached, 2 machines},{8B first load, 2 machines},{3B first load, 2 machines}},
  ytick=data,
  y tick label style={font=\tiny},
  nodes near coords,
  nodes near coords style={font=\tiny,anchor=west,xshift=1pt},
  enlarge y limits=0.12,
  point meta=explicit symbolic,
  ]
\addplot+[fill=serA!55,draw=serA] coordinates {
  (1,{3B rebuild, local only})      [1 s]
  (33,{3B rebuild, 2 machines})     [33 s]
  (48,{3B cached, 3 machines})      [48 s]
  (240,{3B cached, 2 machines})     [240 s]
  (730,{8B first load, 2 machines}) [730 s]
  (990,{3B first load, 2 machines}) [{\textasciitilde}990 s]
};
\end{axis}
\end{tikzpicture}
\caption{Measured time to a ready model. Each bar is a single observation, not
an average. The two first-load bars include transferring the remote worker's
share of the weights over Wi-Fi on a weak link (\SI{-66}{dBm} to
\SI{-68}{dBm}); once the worker's local tensor cache is populated, the same 3B
model reloads in a small fraction of that time. The cached three-machine load
was measured on the stronger \SI{-56}{dBm} link.}
\label{fig:load}
\end{figure}

The practical consequence is that the first load of a new model on a new worker
is slow and should be done before any demonstration, while subsequent loads and
all recoveries are short. Reported cached 3B loads across sessions were
\SI{240}{\second}, \SI{137}{\second}, \SI{60}{\second} and \SI{48}{\second},
decreasing as link conditions improved; we report the range rather than
selecting the most favourable figure.
